% Generated draft: requires mathematical and source review.
\documentclass[11pt]{article}
\usepackage[margin=25mm]{geometry}
\usepackage{fontspec}
\setmainfont{Latin Modern Roman}
\usepackage{amsmath,amssymb,mathtools,unicode-math}
\setmathfont{Latin Modern Math}
\usepackage{enumitem,needspace,xurl,hyperref}
\setlength{\parindent}{0pt}
\setlength{\parskip}{5pt}
\setlength{\emergencystretch}{3em}
\newcommand{\N}{\mathbb{N}}
\newcommand{\Z}{\mathbb{Z}}
\newcommand{\Q}{\mathbb{Q}}
\newcommand{\R}{\mathbb{R}}
\newcommand{\C}{\mathbb{C}}
\newcommand{\F}{\mathbb{F}}
\newcommand{\A}{\mathbb{A}}
\newcommand{\PP}{\mathbb{P}}
\newcommand{\E}{\mathbb{E}}
\newcommand{\eps}{\varepsilon}
\newcommand{\dd}{\,\mathrm{d}}
\begin{document}
\section*{UnsolvedMath ID: 700736}
\textbf{Source problem code:} catalogue entry
\subsection{Definitions and mathematical statement}
For a prime \(p\) and a finite \(p\)-group \(X\), define its sectional rank by
\[
r(X)=\max\bigl\{d\geq 0:\ H/K\cong (\Z/p\Z)^d,\ K\trianglelefteq H\leq X\bigr\}.
\]
Equivalently, this is the maximum, over subgroups of \(X\), of their minimum numbers of generators. In particular, \(r(1)=0\).

A uniform formulation of the question attributed to L. S. Kazarin is: does there exist a real constant \(C\geq 0\) such that, for every prime \(p\), every finite \(p\)-group \(G\), and every pair of subgroups \(A,B\leq G\) satisfying
\[
G=AB=\{ab:a\in A,\ b\in B\},
\]
one has
\[
r(G)\leq C\bigl(r(A)+r(B)\bigr)?
\]
Neither normality of the factors nor trivial intersection is assumed. The displayed bound expresses the existence of a linear bound in the two ranks using a common coefficient. The source does not explicitly specify whether the coefficients must be independent of \(p\); the formulation above requires that uniformity. If dependence on \(p\) is intended, the quantifiers instead become: for every prime \(p\), there exists \(C_p\geq 0\) satisfying the same bound for all such \(G,A,B\).
\subsection{Short English statement}
Kazarin asks whether the sectional rank of a finite \(p\)-group that is a product of two subgroups is bounded linearly by their sectional ranks. The supplied record labels the problem open; the intended uniformity in the prime \(p\) requires confirmation.
\subsection{Sources}
\begin{itemize}
\item Supplied local UnsolvedMath registry record: permanent ID \(700736\), exact source code \texttt{catalogue entry}, title \emph{Kourovka Notebook Problem 11.44}; upstream revision \texttt{d371ee917ac9647fbb06cd6fc529deec6f3dda82}. The record was updated on 2026-10-16 and supplies the status \emph{open}; this status has not been independently verified.
\item Supplied source attribution: \emph{The Kourovka Notebook (Unsolved Problems in Group Theory)}, 11th issue (1990), Problem 11.44, proposed by L. S. Kazarin. The background also identifies the 21st issue edition (2026), arXiv:\texttt{1401.0300v48}. The collection description names E. I. Khukhro and V. D. Mazurov as editors.
\end{itemize}
\end{document}
